% O011-BL — modul asli, bukan terjemahan atau adaptasi dari teks pembanding.
% Judul: Jembatan Grup Lie dan Aljabar Lie
% Lisensi modul: CC BY-SA 4.0.
% Provenans: disusun oleh OpenAI Codex gpt-5.6-sol, Ultra, atas arahan pengguna.
% Prasyarat internal: Unit 7, 9, 10, 11, dan 28 edisi ini.

\setcounter{section}{30}
\section{Jembatan Grup Lie dan Aljabar Lie}
\label{o011-bridge-lie}

\subsection{Tujuan dan hubungan dengan materi sebelumnya}

Modul ini menghubungkan bahasa manifold dengan simetri kontinu. Unit 7
menyediakan manifold mulus, Unit 9 ruang singgung, Unit 10 pemetaan mulus dan
medan vektor, Unit 11 produk manifold, sedangkan Unit 28 menyediakan braket
Lie medan vektor. Semua hasil di bawah dibangun dari perangkat itu. Tujuan
utamanya ialah memahami bagaimana operasi grup yang mulus menghasilkan suatu
aljabar linear pada ruang singgung di unsur identitas, bagaimana eksponensial
menghubungkan kedua struktur tersebut, dan bagaimana aksi grup menghasilkan
orbit serta stabilisator.

Sepanjang modul, semua manifold dianggap berdimensi hingga, Hausdorff, dan
memiliki basis terhitung. Kata ``mulus'' berarti terdiferensialkan tak hingga.
Jika $f\colon M\to N$ mulus dan $p\in M$, diferensialnya ditulis
$T_pf\colon T_pM\to T_{f(p)}N$.

\subsection{Grup Lie dan contoh matriks}

\begin{definition}\label{o011-bl-def-lie-group}
Suatu \emph{grup Lie} ialah sebuah manifold mulus $G$ yang sekaligus merupakan
grup, sedemikian sehingga pemetaan perkalian dan invers
\[
  m\colon G\times G\longrightarrow G,\qquad m(g,h)=gh,
  \qquad
  \iota\colon G\longrightarrow G,\qquad \iota(g)=g^{-1},
\]
bersifat mulus. Unsur identitasnya ditulis $e$. Homomorfisme grup Lie ialah
homomorfisme grup yang juga merupakan pemetaan mulus.
\end{definition}

Syarat tersebut bukan sekadar hiasan topologis. Ia memungkinkan kita
mendiferensialkan identitas grup. Sebagai contoh, dari $gg^{-1}=e$ kita akan
memperoleh rumus diferensial untuk invers; dari keasosiatifan kita akan
memperoleh medan vektor invarian dan braket pada $T_eG$.

\begin{beispiel}\label{o011-bl-ex-vector-groups}
Ruang vektor nyata $\mathbb R^n$ dengan penjumlahan merupakan grup Lie. Unsur
identitasnya $0$, inversnya $x\mapsto -x$, dan kedua operasi itu polinomial.
Torus
\[
  \mathbb T^n=(S^1)^n
\]
dengan perkalian komponen demi komponen juga merupakan grup Lie. Pemetaan
\[
  \mathbb R^n\longrightarrow\mathbb T^n,
  \qquad (t_1,\ldots,t_n)\longmapsto
  (e^{it_1},\ldots,e^{it_n})
\]
ialah homomorfisme grup Lie surjektif dengan kernel $(2\pi\mathbb Z)^n$.
\end{beispiel}

\begin{beispiel}\label{o011-bl-ex-general-linear}
Di dalam ruang semua matriks nyata $M_n(\mathbb R)\cong\mathbb R^{n^2}$,
himpunan
\[
  \mathrm{GL}_n(\mathbb R)=\{A\in M_n(\mathbb R):\det A\ne0\}
\]
terbuka karena determinan kontinu. Jadi ia merupakan manifold. Perkalian
matriks bersifat polinomial, sedangkan
\[
  A^{-1}=\frac{1}{\det A}\operatorname{adj}(A)
\]
mulus pada himpunan tempat $\det A\ne0$. Dengan demikian
$\mathrm{GL}_n(\mathbb R)$ ialah grup Lie.
\end{beispiel}

\begin{beispiel}\label{o011-bl-ex-matrix-subgroups}
Contoh-contoh penting berikut merupakan subgrup tertutup
$\mathrm{GL}_n(\mathbb R)$ dan, dengan struktur manifold terinduksi yang
sesuai, merupakan grup Lie:
\[
\begin{aligned}
 \mathrm{SL}_n(\mathbb R)&=\{A:\det A=1\},\\
 \mathrm O(n)&=\{A:A^{\mathsf T}A=I\},\\
 \mathrm{SO}(n)&=\{A:A^{\mathsf T}A=I,\ \det A=1\}.
\end{aligned}
\]
Untuk $\mathrm{SL}_n$, nilai $1$ merupakan nilai regular determinan karena
$D(\det)_A(H)=\det(A)\operatorname{tr}(A^{-1}H)$ tidak nol sebagai fungsional.
Untuk $\mathrm O(n)$, persamaan $A^{\mathsf T}A=I$ dapat diperlakukan dengan
teorema nilai regular ke ruang matriks simetris. Kita tidak memerlukan teorema
umum bahwa setiap subgrup tertutup dari grup Lie adalah submanifold, meskipun
teorema itu memang benar.
\end{beispiel}

\begin{bemerkung}\label{o011-bl-rem-components}
$\mathrm{GL}_n(\mathbb R)$ mempunyai dua komponen terhubung, dibedakan oleh
tanda determinan. $\mathrm O(n)$ juga mempunyai dua komponen, sedangkan
$\mathrm{SO}(n)$ adalah komponen identitasnya. Aljabar Lie hanya merekam
geometri infinitesimal di sekitar $e$; karena itu ia tidak dengan sendirinya
merekam semua komponen terhubung suatu grup.
\end{bemerkung}

\subsection{Translasi kiri dan medan vektor invarian}

Untuk $g\in G$, definisikan translasi kiri dan kanan
\[
  L_g(h)=gh,
  \qquad
  R_g(h)=hg.
\]
Keduanya difeomorfisme, dengan invers masing-masing $L_{g^{-1}}$ dan
$R_{g^{-1}}$.

\begin{definition}\label{o011-bl-def-left-invariant}
Suatu medan vektor mulus $X$ pada $G$ disebut \emph{invarian-kiri} apabila
\[
  T_hL_g(X_h)=X_{gh}
\]
untuk semua $g,h\in G$.
\end{definition}

\begin{Proposition}\label{o011-bl-prop-left-invariant-correspondence}
Untuk setiap $v\in T_eG$ terdapat tepat satu medan vektor invarian-kiri
$X^v$, yaitu
\[
  X^v_g=T_eL_g(v).
\]
Pemetaan $v\mapsto X^v$ merupakan isomorfisme linear dari $T_eG$ ke ruang
medan vektor invarian-kiri.
\end{Proposition}

\begin{proof}
Keasosiatifan memberi $L_g\circ L_h=L_{gh}$. Maka
\[
 T_hL_g(X^v_h)
 =T_hL_g\bigl(T_eL_h(v)\bigr)
 =T_e(L_g\circ L_h)(v)
 =T_eL_{gh}(v)=X^v_{gh}.
\]
Sebaliknya, jika $X$ invarian-kiri, ambil $h=e$ untuk memperoleh
$X_g=T_eL_g(X_e)$. Jadi $X$ ditentukan secara unik oleh nilainya di $e$.
Linearitas jelas dari linearitas diferensial.
\end{proof}

\begin{Proposition}\label{o011-bl-prop-invariant-bracket}
Jika $X$ dan $Y$ invarian-kiri, maka braket medan vektor $[X,Y]$ juga
invarian-kiri.
\end{Proposition}

\begin{proof}
Untuk setiap difeomorfisme $F$, sifat natural braket menyatakan
$F_*[X,Y]=[F_*X,F_*Y]$. Terapkan ini pada $F=L_g$. Karena
$(L_g)_*X=X$ dan $(L_g)_*Y=Y$, diperoleh $(L_g)_*[X,Y]=[X,Y]$.
\end{proof}

\subsection{Aljabar Lie di unsur identitas}

\begin{definition}\label{o011-bl-def-lie-algebra}
\emph{Aljabar Lie} suatu grup Lie $G$ ialah ruang vektor
\[
  \mathfrak g=T_eG
\]
dengan operasi bilinear
\[
  [v,w]_{\mathfrak g}=[X^v,X^w]_e.
\]
Jika konteksnya jelas, subskrip $\mathfrak g$ dihilangkan.
\end{definition}

Proposisi sebelumnya menjamin bahwa $[X^v,X^w]$ kembali invarian-kiri,
sehingga nilainya di $e$ menentukan seluruh medan. Sifat braket medan vektor
langsung memberi
\[
 [v,w]=-[w,v],
 \qquad
 [u,[v,w]]+[v,[w,u]]+[w,[u,v]]=0.
\]
Identitas kedua disebut identitas Jacobi.

\begin{Proposition}\label{o011-bl-prop-homomorphism-differential}
Jika $F\colon G\to H$ homomorfisme grup Lie, maka
\[
  dF_e=T_eF\colon\mathfrak g\longrightarrow\mathfrak h
\]
mempertahankan braket Lie.
\end{Proposition}

\begin{proof}
Identitas $F\circ L_g=L_{F(g)}\circ F$ menunjukkan bahwa medan
invarian-kiri $X^v$ dan $X^{dF_e(v)}$ saling terkait oleh $F$. Naturalitas
braket untuk medan yang saling terkait menghasilkan
\[
  dF_e([v,w])=[dF_e(v),dF_e(w)].
\]
\end{proof}

\begin{beispiel}\label{o011-bl-ex-matrix-algebras}
Untuk grup matriks, kita mengidentifikasi ruang singgung di $I$ dengan
subruang matriks. Hasilnya ialah
\[
\begin{aligned}
 \mathfrak{gl}_n(\mathbb R)&=M_n(\mathbb R),\\
 \mathfrak{sl}_n(\mathbb R)&=\{X:\operatorname{tr}X=0\},\\
 \mathfrak o(n)=\mathfrak{so}(n)&=\{X:X^{\mathsf T}+X=0\}.
\end{aligned}
\]
Braketnya adalah komutator
\[
  [X,Y]=XY-YX.
\]
Sebagai contoh, diferensiasikan $A(t)^{\mathsf T}A(t)=I$ pada $t=0$ dengan
$A(0)=I$ untuk memperoleh $A'(0)^{\mathsf T}+A'(0)=0$.
\end{beispiel}

\begin{bemerkung}\label{o011-bl-rem-sign}
Konvensi braket di atas memakai medan invarian-kiri. Jika seseorang
mengidentifikasi $T_eG$ dengan medan invarian-kanan tanpa mengubah definisi,
tanda braket akan berbalik. Menyatakan konvensi ini mencegah kesalahan tanda
pada rumus adjoint.
\end{bemerkung}

\subsection{Subgrup satu-parameter dan pemetaan eksponensial}

\begin{definition}\label{o011-bl-def-one-parameter}
Subgrup satu-parameter dalam $G$ ialah homomorfisme grup Lie
\[
  \gamma\colon(\mathbb R,+)\longrightarrow G.
\]
\end{definition}

Jika $\gamma$ subgrup satu-parameter, maka
\[
  \gamma(t+s)=\gamma(t)\gamma(s)
\]
dan, setelah mendiferensialkan terhadap $s$ pada $0$,
\[
  \gamma'(t)=T_eL_{\gamma(t)}\bigl(\gamma'(0)\bigr).
\]
Jadi $\gamma$ adalah kurva integral medan invarian-kiri yang ditentukan oleh
$v=\gamma'(0)$.

\begin{Theorem}\label{o011-bl-thm-one-parameter-existence}
Untuk setiap $v\in\mathfrak g$ terdapat tepat satu subgrup satu-parameter
$\gamma_v$ dengan $\gamma_v'(0)=v$.
\end{Theorem}

\begin{proof}
Ambil kurva integral maksimal $\gamma$ dari medan invarian-kiri $X^v$ dengan
$\gamma(0)=e$. Untuk $s$ tetap, kurva $t\mapsto\gamma(s)\gamma(t)$ dan
$t\mapsto\gamma(s+t)$ memenuhi persamaan diferensial yang sama serta memiliki
nilai awal yang sama pada $t=0$. Ketunggalan solusi memberi
$\gamma(s+t)=\gamma(s)\gamma(t)$ selama kedua ruas terdefinisi.

Identitas ini memperpanjang kurva integral lokal langkah demi langkah ke
seluruh $\mathbb R$: pilih $\varepsilon>0$ dengan solusi pada
$(-\varepsilon,\varepsilon)$, lalu definisikan nilai pada selang yang lebih
panjang dengan hasil kali potongan-potongan kecil. Identitas lokal dan
ketunggalan menjamin bahwa definisi itu konsisten. Kurva global yang diperoleh
adalah homomorfisme. Ketunggalan subgrup satu-parameter juga mengikuti dari
ketunggalan kurva integral.
\end{proof}

\begin{definition}\label{o011-bl-def-exponential}
Pemetaan eksponensial grup Lie adalah
\[
  \exp_G\colon\mathfrak g\longrightarrow G,
  \qquad \exp_G(v)=\gamma_v(1).
\]
\end{definition}

Dari definisi dan penskalaan parameter diperoleh
\[
  \gamma_v(t)=\exp_G(tv),
  \qquad
  \exp_G((s+t)v)=\exp_G(sv)\exp_G(tv).
\]

\begin{Proposition}\label{o011-bl-prop-exp-local}
Pemetaan $\exp_G$ mulus, $\exp_G(0)=e$, dan diferensialnya di $0$ adalah
identitas pada $\mathfrak g$. Karena itu $\exp_G$ merupakan difeomorfisme dari
suatu lingkungan $0$ ke suatu lingkungan $e$.
\end{Proposition}

\begin{proof}
Kebergantungan mulus solusi persamaan diferensial pada nilai awal dan parameter
memberi kemulusan $\exp_G$. Kurva $t\mapsto\exp_G(tv)$ mempunyai kecepatan $v$
pada $0$, sehingga
\[
  T_0\exp_G(v)=\left.\frac{d}{dt}\right|_{0}\exp_G(tv)=v.
\]
Teorema fungsi invers menyelesaikan pernyataan terakhir.
\end{proof}

\begin{beispiel}\label{o011-bl-ex-matrix-exp}
Untuk $G=\mathrm{GL}_n(\mathbb R)$,
\[
  \exp_G(X)=e^X=\sum_{k=0}^{\infty}\frac{X^k}{k!}.
\]
Deret ini konvergen mutlak dalam setiap norma matriks. Turunan
$t\mapsto e^{tX}$ ialah $Xe^{tX}=e^{tX}X$, sehingga kurva tersebut merupakan
subgrup satu-parameter dengan kecepatan awal $X$. Jika $X$ dan $Y$ komutatif,
$e^{X+Y}=e^Xe^Y$; tanpa hipotesis komutatif, identitas ini umumnya salah.
\end{beispiel}

\begin{Proposition}\label{o011-bl-prop-exp-natural}
Untuk homomorfisme grup Lie $F\colon G\to H$ berlaku
\[
  F(\exp_G v)=\exp_H(dF_e v).
\]
\end{Proposition}

\begin{proof}
Kedua ruas, setelah $v$ diganti $tv$, merupakan subgrup satu-parameter di $H$
dengan kecepatan awal $dF_e v$. Gunakan ketunggalan pada Teorema
\ref{o011-bl-thm-one-parameter-existence} dan ambil $t=1$.
\end{proof}

\subsection{Konjugasi dan representasi adjoint}

Untuk $g\in G$, konjugasi oleh $g$ adalah difeomorfisme
\[
  C_g\colon G\longrightarrow G,
  \qquad C_g(h)=ghg^{-1}.
\]
Ia mempertahankan $e$.

\begin{definition}\label{o011-bl-def-adjoint}
Pemetaan adjoint pada tingkat grup adalah
\[
  \operatorname{Ad}\colon G\longrightarrow\mathrm{GL}(\mathfrak g),
  \qquad \operatorname{Ad}_g=T_eC_g.
\]
Diferensialnya di identitas adalah
\[
  \operatorname{ad}=T_e\operatorname{Ad}\colon
  \mathfrak g\longrightarrow\mathfrak{gl}(\mathfrak g).
\]
\end{definition}

Karena $C_{gh}=C_g\circ C_h$, pemetaan $\operatorname{Ad}$ merupakan
homomorfisme grup Lie. Untuk grup matriks,
\[
  \operatorname{Ad}_g(X)=gXg^{-1}.
\]

\begin{Theorem}\label{o011-bl-thm-ad-bracket}
Untuk $v,w\in\mathfrak g$ berlaku
\[
  \operatorname{ad}_v(w)=[v,w].
\]
Selain itu,
\[
  \operatorname{Ad}_g[v,w]
  =[\operatorname{Ad}_gv,\operatorname{Ad}_gw]
\]
dan
\[
  \operatorname{Ad}_{\exp(tv)}
  =\exp\bigl(t\operatorname{ad}_v\bigr).
\]
\end{Theorem}

\begin{proof}
Konjugasi membawa medan invarian-kiri $X^w$ ke medan invarian-kiri yang
bernilai $\operatorname{Ad}_g w$ di identitas. Naturalitas braket memberi
identitas kedua. Untuk identitas pertama, diferensiasikan keluarga
$C_{\exp(tv)}$ pada $t=0$; turunan aksi keluarga difeomorfisme ini pada medan
invarian-kiri adalah braket dengan $X^v$. Evaluasi di $e$ memberi
$\operatorname{ad}_v(w)=[v,w]$.

Terakhir, $t\mapsto\operatorname{Ad}_{\exp(tv)}$ ialah subgrup satu-parameter
di $\mathrm{GL}(\mathfrak g)$ dengan kecepatan awal
$\operatorname{ad}_v$. Contoh \ref{o011-bl-ex-matrix-exp} pada ruang
$\mathfrak g$ memberi identitas ketiga.
\end{proof}

\subsection{Aksi mulus, orbit, dan stabilisator}

\begin{definition}\label{o011-bl-def-action}
Aksi kiri mulus grup Lie $G$ pada manifold $M$ adalah pemetaan mulus
\[
  \Phi\colon G\times M\longrightarrow M,
  \qquad (g,p)\longmapsto g\cdot p,
\]
dengan $e\cdot p=p$ dan $(gh)\cdot p=g\cdot(h\cdot p)$.
Orbit dan stabilisator titik $p$ adalah
\[
  G\cdot p=\{g\cdot p:g\in G\},
  \qquad
  G_p=\{g\in G:g\cdot p=p\}.
\]
\end{definition}

Stabilisator $G_p$ adalah subgrup tertutup, sebab ia adalah prapeta
$\{p\}$ di bawah pemetaan kontinu $g\mapsto g\cdot p$. Pemetaan orbit
\[
  \Phi_p\colon G\longrightarrow M,
  \qquad \Phi_p(g)=g\cdot p
\]
memenuhi $\Phi_p(gh)=g\cdot\Phi_p(h)$.

\begin{definition}\label{o011-bl-def-fundamental-field}
Untuk $v\in\mathfrak g$, medan vektor fundamental yang dihasilkan oleh $v$
adalah
\[
  v_M(p)=\left.\frac{d}{dt}\right|_{0}\exp(tv)\cdot p
  =T_e\Phi_p(v).
\]
\end{definition}

Dengan konvensi aksi kiri dan medan invarian-kiri yang dipakai di sini,
pemetaan $v\mapsto v_M$ merupakan antihomomorfisme aljabar Lie:
\[
  [v_M,w_M]=-[v,w]_M.
\]
Tanda minus berasal dari fakta bahwa $p\mapsto\exp(tv)\cdot p$ membentuk
aksi kiri, sedangkan medan fundamentalnya terkait secara alami dengan medan
invarian-kanan pada $G$. Mengubah salah satu konvensi akan mengubah tanda.

\begin{Proposition}\label{o011-bl-prop-orbit-tangent}
Untuk setiap $p\in M$,
\[
  \operatorname{im}(T_e\Phi_p)
  =\{v_M(p):v\in\mathfrak g\}
\]
adalah ruang arah tangensial orbit di $p$. Kernel
\[
  \mathfrak g_p=\ker(T_e\Phi_p)
\]
adalah aljabar Lie stabilisator. Jika orbit mempunyai struktur manifold
homogen yang lazim, maka
\[
  T_p(G\cdot p)\cong\mathfrak g/\mathfrak g_p,
  \qquad
  \dim(G\cdot p)=\dim G-\dim G_p.
\]
\end{Proposition}

\begin{proof}
Identitas pertama adalah definisi diferensial pemetaan orbit. Kurva
$\exp(tv)$ terletak di $G_p$ tepat secara infinitesimal ketika
$v_M(p)=0$, sehingga kernelnya adalah ruang singgung stabilisator di $e$.
Teorema isomorfisme linear memberi
$\mathfrak g/\ker(T_e\Phi_p)\cong\operatorname{im}(T_e\Phi_p)$.
Apabila orbit direalisasikan sebagai manifold homogen $G/G_p$, diferensial
pemetaan hasil bagi mengidentifikasi ruang terakhir dengan $T_p(G\cdot p)$.
Rumus dimensi mengikuti.
\end{proof}

\begin{beispiel}\label{o011-bl-ex-sphere-action}
$\mathrm{SO}(n)$ bertindak pada $S^{n-1}$ melalui perkalian matriks. Aksi ini
transitif. Stabilisator $e_n$ terdiri atas matriks berbentuk
\[
  \begin{pmatrix} A&0\\0&1\end{pmatrix},
  \qquad A\in\mathrm{SO}(n-1),
\]
sehingga $S^{n-1}\cong\mathrm{SO}(n)/\mathrm{SO}(n-1)$. Rumus dimensi memberi
\[
 \frac{n(n-1)}2-\frac{(n-1)(n-2)}2=n-1,
\]
sesuai dengan dimensi bola satuan.
\end{beispiel}

\subsection{Ringkasan struktural}

Data global dan infinitesimal yang dibangun di atas tersusun sebagai berikut:
\[
\begin{array}{ccl}
G&\longmapsto&\mathfrak g=T_eG,\\
F\colon G\to H&\longmapsto&dF_e\colon\mathfrak g\to\mathfrak h,\\
v\in\mathfrak g&\longmapsto&X^v\text{ dan }\exp(tv),\\
g\in G&\longmapsto&\operatorname{Ad}_g,\\
v\in\mathfrak g&\longmapsto&\operatorname{ad}_v=[v,\cdot],\\
G\curvearrowright M&\longmapsto&v_M(p)=T_e\Phi_p(v).
\end{array}
\]
Eksponensial memberi koordinat lokal di sekitar identitas, tetapi tidak harus
surjektif secara global dan tidak mengubah persoalan global menjadi persoalan
linear. Orbit mengubah simetri menjadi geometri submanifold, sedangkan
stabilisator mengukur simetri yang tidak menggerakkan suatu titik.
